\begin{abstract}
Deciding when to remove a jet engine from service trades the cost of an
in-service failure against the value of the life left unused. Deep-learning
models for remaining-useful-life (RUL) prediction report steadily lower error
on the NASA C-MAPSS benchmark, but they output point estimates without
calibrated confidence, and they are evaluated on prediction error rather than
on the maintenance decisions they would drive. This paper presents
\emph{Twin Turbo}, a framework that combines a lightweight edge model with a
per-engine digital twin and makes autonomous maintenance decisions from
calibrated lower bounds. The edge model is a gradient-boosted tree ensemble
over 30-cycle sensor windows, wrapped in Mondrian conformal prediction so that
its intervals carry a distribution-free coverage guarantee. The digital twin is
a particle filter over an exponential health-index degradation model: a
population prior learned from run-to-failure engines is updated with the
observations of one specific engine, yielding a full RUL distribution. The edge
model runs every cycle; the twin is synchronized only when one of three
triggers fires, and the policy acts on the more cautious of the two calibrated
bounds. On the official C-MAPSS test sets the twin reaches an RMSE of
10.33 on FD002, and its 90\% intervals achieve 86--94\% empirical coverage,
against 43--54\% for Monte-Carlo dropout on an LSTM baseline. Replaying 709
held-out engines from first cycle to failure through five-fold
cross-validation, the hybrid policy lets no engine fail, lowers the long-run
maintenance cost rate by 34--47\% relative to optimal age replacement, consults
the twin on only 21--28\% of cycles, and is the only policy evaluated that
remains failure-free without any safety buffer on all four datasets.
\end{abstract}

\begin{IEEEkeywords}
prognostics and health management, remaining useful life, digital twin,
particle filter, conformal prediction, edge computing, predictive maintenance,
C-MAPSS
\end{IEEEkeywords}
